\documentclass[11pt]{article}
\usepackage{amsmath,amssymb,amsthm,hyperref}
\newtheorem{proposition}{Proposition}
\newcommand{\per}{\operatorname{per}}
\title{An exact obstruction to gap-monotonicity induction for Marcus's conjecture}
\author{Research notes for independent verification; no priority claim}
\date{6 November 2026}
\begin{document}
\maketitle
\noindent\textbf{Status.} This does not prove or disprove Marcus's
block-permanent conjecture. It gives an exact failure of a stronger
inductive inequality, using split quadratic forms and positive definite
diagonal blocks. The full conjecture remains the goal.

For homogeneous polynomials of equal positive degree, use the Fock inner
product, conjugate-linear in its first argument, with
$\langle z^\alpha,z^\beta\rangle=\mathbf1_{\alpha=\beta}\alpha!$, and put
\[
 F_m(p_1,\ldots,p_m)=\|p_1\cdots p_m\|^2
             -\per[\langle p_i,p_j\rangle]_{i,j=1}^m.
\]
When every $p_i$ is a product of $k$ linear forms, this is the gap in
Marcus's conjecture for the Gram matrix of their $mk$ coefficient vectors.

\paragraph{The attempted induction.}
A sufficient inductive bound would have been
\begin{equation}\label{bad}
 F_m(p_1,\ldots,p_m)\ge
 F_{m-1}(p_1,\ldots,p_{m-1})\|p_m\|^2.
\end{equation}
Since $F_1=0$, this would prove the inequality for all $m$.
It fails already at $m=3,k=2$, including with positive definite diagonal
blocks, as shown below.

The exact operator behind the attempt can still be useful. Set
$f=\prod_{i<m}p_i$, $H=[\langle p_i,p_j\rangle]_{i,j<m}$,
$q=\per H$, and $c_{ij}=\per H^{(i|j)}$, where the indicated row and
column are removed and the empty permanent is one. For a degree-$k$ form $g$,
expansion according to how the last row and column of the permanent are
matched gives
\[
 \per\operatorname{Gram}(p_1,\ldots,p_{m-1},g)
 =q\|g\|^2+\sum_{i,j<m}c_{ij}
       \langle p_i,g\rangle\langle g,p_j\rangle.
\]
Consequently, with
\[
 \langle g,Tg\rangle=
 \sum_{i,j<m}c_{ij}\langle p_i,g\rangle\langle g,p_j\rangle,
 \qquad K=M_f^*M_f-\|f\|^2I,
\]
where $M_f$ is multiplication by $f$, the exact remainder is
\begin{equation}\label{remainder}
 F_m=F_{m-1}\|g\|^2+\langle g,(K-T)g\rangle.
\end{equation}
The normal-ordering identity
\[
 \|fg\|^2=\sum_\alpha\frac1{\alpha!}
       \|(\partial^\alpha f)^*(\partial)g\|^2
\]
shows $K\succeq0$: the terms $|\alpha|=\deg f$ sum to
$\|f\|^2\|g\|^2$ and all other terms are nonnegative.
However, $K\succeq T$ is false. Dropping the scalar slack
$F_{m-1}I$ from the full operator comparison loses essential information.
Nonnegativity of $F_m$ for every homogeneous degree-$k$ form $g$ would
be equivalent to $K+F_{m-1}I\succeq T$. This operator inequality is not
proved here. Marcus's original statement requires the quadratic-form
inequality only on split $g$.

\begin{proposition}[Exact family of failures]
Let $A,C>0$ and $B=\sqrt{2AC}$ be real. Set
\[
 p_1=Ax^2+Bxy+Cy^2,\quad
 p_2=Ax^2-Bxy+Cy^2,\quad
 g=cx^2+dy^2
\]
with real $c,d$. Then
\begin{align}\label{family}
 F_3(p_1,p_2,g)-F_2(p_1,p_2)\|g\|^2
 =32\bigl[21A^4c^2+21C^4d^2
             -(A^2+C^2)(Ac+Cd)^2\bigr].
\end{align}
\end{proposition}
\begin{proof}
The imposed relation cancels the mixed term in
$f=p_1p_2=A^2x^4+C^2y^4$, giving
\[
 \|f\|^2=24(A^4+C^4),\qquad
 \|fg\|^2=(720A^4+48C^4)c^2+(48A^4+720C^4)d^2.
\]
Writing $n=2(A^2+AC+C^2)$ and $s=2(A^2-AC+C^2)$, the first
two forms have Gram matrix $\left(\begin{smallmatrix}n&s\\s&n\end{smallmatrix}\right)$.
Both their inner products with $g$ equal $t=2(Ac+Cd)$.
The third norm is $\|g\|^2=2(c^2+d^2)$. Therefore
\[
 \per\operatorname{Gram}(p_1,p_2,g)
 =\|g\|^2(n^2+s^2)+2(n+s)t^2.
\]
Subtracting $F_2\|g\|^2$ cancels the first permanent term and yields
\eqref{family}.
\end{proof}

\paragraph{A Gaussian-integer witness with nonsingular diagonal blocks.}
Choose $A=8,C=1,B=4,c=1,d=36$. The six coefficient vectors are
\[
 V=\begin{pmatrix}
 2+2i&1\\2-2i&1\\-2+2i&1\\-2-2i&1\\1&6i\\1&-6i
 \end{pmatrix}.
\]
Each consecutive pair gives one quadratic form. Let
$\mathcal A_{ij}=\langle v_i,v_j\rangle$ and partition this positive
semidefinite Gram matrix into three pairs. The determinants of its three
diagonal blocks are $16,16,144$, so they are all positive definite.
The matrix of block permanents is
\[
 G=\begin{pmatrix}146&114&88\\114&146&88\\88&88&2594\end{pmatrix}.
\]
Exact values are
\[
 \per\mathcal A=258686256,\quad \per G=93032208,\quad
 F_3=165654048,
\]
whereas
\[
 F_2=64016,\quad\|g\|^2=2594,\quad
 F_2\|g\|^2=166057504.
\]
Thus the proposed increment is $-403456$, even though $F_3$ is positive.
This is not a counterexample to Marcus's conjecture.

The script \texttt{marcus-induction-obstruction.mjs} verifies these
quantities with Gaussian-integer arithmetic. It compares the original Gram
permanent, independent permutation enumeration, and polynomial convolution
with factorial norms. It also checks \eqref{family} at $100$ integer-factor
specializations $A=2t^2,C=1,B=2t,c=1,d=s^2$, $1\le t\le5$,
$1\le s\le20$. The symbolic derivation proves the family; the checks
do not extend its scope.

\paragraph{Implication for the full problem.}
An induction using \eqref{remainder} must control the negative part of
$K-T$ using the previous gap, or use a different decomposition. Positivity
of $K$ alone, and nonnegativity of the previous gap alone, do not accomplish
this. No adequate uniform control is proved here.
\end{document}
